% Paste this environment into the IEEEtran source.
% The preamble must already \input the sibling iscafig.tex.
\begin{figure*}[t]
\centering
\begin{tikzpicture}[iscafig]
  \begin{scope}
    \node[lbl, anchor=west] at (0,0.40) {\textbf{Instance}};
    \rec{0}{+0}{fldhot}{5.15cm}{kind $= 1$}
    \rec{-0.55}{+4}{fldmem}{5.15cm}{type}
    \rec{-1.10}{+8}{fldmem}{5.15cm}{\texttt{\_\_dict\_\_}, or null}
    \node[lbl, anchor=west] at (0,-1.85) {\textbf{Exception}};
    \rec{-2.25}{+0}{fldhot}{5.15cm}{kind $= 3$}
    \rec{-2.80}{+4}{fldmem}{5.15cm}{type}
    \rec{-3.35}{+8}{fldmem}{5.15cm}{\texttt{args} tuple}
    \rec{-3.90}{+12}{fldmem}{5.15cm}{traceback}
    \rec{-4.45}{+16}{fldmem}{5.15cm}{\texttt{\_\_dict\_\_}, or null}
  \end{scope}

  \begin{scope}[xshift=6.15cm]
    \node[lbl, anchor=west] at (0,0.40) {\textbf{Type}, nine words};
    \rec{0}{+0}{fldhot}{5.15cm}{kind $= 2$}
    \rec{-0.55}{+4}{fldmem}{5.15cm}{metaclass}
    \rec{-1.10}{+8}{fldmem}{5.15cm}{namespace dict}
    \rec{-1.65}{+12}{fldmem}{5.15cm}{bases tuple}
    \rec{-2.20}{+16}{fldmem}{5.15cm}{name}
    \rec{-2.75}{+20}{fldmem}{5.15cm}{MRO tuple}
    \rec{-3.30}{+24}{fldimm}{5.15cm}{reserved}
    \rec{-3.85}{+28}{fldimm}{5.15cm}{flags}
    \rec{-4.40}{+32}{fldmem}{5.15cm}{\texttt{\_\_new\_\_}, or null}
  \end{scope}

  \begin{scope}[xshift=12.3cm]
    \node[lbl, anchor=west] at (0,0.40) {\textbf{Function}, eight words};
    \rec{0}{+0}{fldhot}{5.15cm}{kind $= 5$}
    \rec{-0.55}{+4}{fldmem}{5.15cm}{type}
    \rec{-1.10}{+8}{fldimm}{5.15cm}{raw code, immediate}
    \rec{-1.65}{+12}{fldmem}{5.15cm}{closure, or null}
    \rec{-2.20}{+16}{fldmem}{5.15cm}{defaults, or null}
    \rec{-2.75}{+20}{fldmem}{5.15cm}{keyword defaults}
    \rec{-3.30}{+24}{fldmem}{5.15cm}{globals}
    \rec{-3.85}{+28}{fldmem}{5.15cm}{root tuple}
    \node[lbl, anchor=west, align=left, text width=5.1cm] at (0,-4.55)
      {A custom \texttt{\_\_hash\_\_} is an entry of the type's namespace dict, not a header slot.};
  \end{scope}
\end{tikzpicture}
\caption{Heap objects under the \texttt{OBJECT} tag.
Word 0 is the kind and word 1 is the type pointer.
Fields start at word 2.
A bound method is four words: kind 7, type, function, and \texttt{self}.
A generator is six words: kind 6, type, function, status, pending exception, and a pointer to its saved frame.
Green boxes are pointers.
The raw-code word is an immediate index into flash, so the collector does not follow it.}
\label{fig:object}
\end{figure*}
